% \documentclass[a4paper,12pt]{article}{{{
\documentclass[a4paper,11pt]{article}
% -------------------------------------------------
% Set up the page margins.
% -------------------------------------------------
% \usepackage[text={6.5in,9.5in},centering]{geometry}
% \setlength{\topmargin}{-0.25in}
\usepackage[margin=1in]{geometry}

% -------------------------------------------------
% Load Packages here
% -------------------------------------------------
%\usepackage{hyperref}
\usepackage{graphicx,url,subfig}
\RequirePackage[OT1]{fontenc}
\RequirePackage{amsthm,amsmath,amssymb,amscd}
\RequirePackage[numbers]{natbib}
\RequirePackage[colorlinks,citecolor=blue,urlcolor=blue]{hyperref}
\usepackage{graphics}
\usepackage{enumerate}
\usepackage{datetime}
\usepackage{etex}
% \usepackage[notcite,notref,color]{showkeys}
\usepackage[normalem]{ulem} % either use this (simple) or
\usepackage{soul} % use this (many fancier options)
\makeatother
\numberwithin{equation}{section}
\allowdisplaybreaks
\usepackage{mathrsfs}

%copied from invariant measure paper
%-----
\usepackage{amssymb, latexsym, amsmath, graphics, fullpage, epsfig, amsthm, relsize, pgf, tikz, amsfonts, makeidx, latexsym, ifthen, calc}
\usepackage[colorlinks,citecolor=blue,urlcolor=blue]{hyperref}
%\usepackage{eucal} this is a different font than \mathcal{•}
\usetikzlibrary{arrows}
\usepackage{mathrsfs}
\usepackage[shortlabels]{enumitem}
\usepackage{tikz}
\usepackage{amsmath}
\usetikzlibrary{arrows}
\usetikzlibrary{shadows.blur}
\usepackage{pgfplots}
\usepackage{mwe} % For dummy images
% \usepackage{subcaption}
\usepackage{multirow}
\usepackage{dcolumn}
\newcolumntype{2}{D{.}{}{2.0}}
\usepackage{multicol}
\numberwithin{equation}{section}
%----
%end copy from invariant measure paper

% -------------------------------------------------
% check references
% -------------------------------------------------
% \usepackage{refcheck}
% \usepackage[notref,notcite]{showkeys}
% \usepackage[notcite]{showkeys}
% \usepackage{showkeys}

% -------------------------------------------------
% Environments here
% -------------------------------------------------
\theoremstyle{plain}
\newtheorem{theorem}{Theorem}[section]
\newtheorem{corollary}[theorem]{Corollary}
\newtheorem{lemma}[theorem]{Lemma}
\newtheorem{proposition}[theorem]{Proposition}
\newtheorem{exercise}{Exercise}
\theoremstyle{definition}
\newtheorem{definition}[theorem]{Definition}
\newtheorem{hypothesis}[theorem]{Hypothesis}
\newtheorem{assumption}[theorem]{Assumption}

\newtheorem{remark}[theorem]{Remark}
\newtheorem{conjecture}[theorem]{Conjecture}
\newtheorem{example}[theorem]{Example}

% -------------------------------------------------
%  Define short hand symbols.
% -------------------------------------------------
\newcommand{\E}{\mathbb{E}}
\newcommand{\W}{\dot{W}}
\newcommand{\B}{\textbf{B}}
\newcommand{\D}{\mathbb{D}}
\newcommand{\ud}{\ensuremath{\mathrm{d} }}
\newcommand{\Ceil}[1]{\left\lceil #1 \right\rceil}
\newcommand{\Floor}[1]{\left\lfloor #1 \right\rfloor}
\newcommand{\sgn}{\text{sgn}}
\newcommand{\Lad}{\text{L}_{\text{ad}}^2}
\newcommand{\SI}[1]{\mathcal{I}\left[#1 \right]}
\newcommand{\SIB}[2]{\mathcal{I}_{#2}\left[#1 \right]}
\newcommand{\Indt}[1]{1_{\left\{#1 \right\} }}
\newcommand{\LadInPrd}[1]{\left\langle #1 \right\rangle_{\text{L}_\text{ad}^2}}
\newcommand{\LadNorm}[1]{\left|\left|  #1 \right|\right|_{\text{L}_\text{ad}^2}}
\newcommand{\Norm}[1]{\left\|  #1   \right\|}
\newcommand{\Ito}{It\^{o} }
\newcommand{\Itos}{It\^{o}'s }
\newcommand{\spt}[1]{\text{supp}\left(#1\right)}
\newcommand{\InPrd}[1]{\left\langle #1 \right\rangle}
\newcommand{\mr}{\textbf{r}}
\newcommand{\Ei}{\text{Ei}}
\newcommand{\arctanh}{\operatorname{arctanh}}
\newcommand{\ind}[1]{\mathbb{I}_{\left\{ {#1} \right\} }}

\DeclareMathOperator{\esssup}{\ensuremath{ess\,sup}}

\newcommand{\steps}[1]{\vskip 0.3cm \textbf{#1}}
\newcommand{\calB}{\mathcal{B}}
\newcommand{\calD}{\mathcal{D}}
\newcommand{\calE}{\mathcal{E}}
\newcommand{\calF}{\mathcal{F}}
\newcommand{\calG}{\mathcal{G}}
\newcommand{\calK}{\mathcal{K}}
\newcommand{\calH}{\mathcal{H}}
\newcommand{\calI}{\mathcal{I}}
\newcommand{\calL}{\mathcal{L}}
\newcommand{\calM}{\mathcal{M}}
\newcommand{\calN}{\mathcal{N}}
\newcommand{\calO}{\mathcal{O}}
\newcommand{\calT}{\mathcal{T}}
\newcommand{\calP}{\mathcal{P}}
\newcommand{\calR}{\mathcal{R}}
\newcommand{\calS}{\mathcal{S}}
\newcommand{\bbC}{\mathbb{C}}
\newcommand{\bbN}{\mathbb{N}}
\newcommand{\bbP}{\mathbb{P}}
\newcommand{\bbZ}{\mathbb{Z}}
\newcommand{\myVec}[1]{\overrightarrow{#1}}
\newcommand{\sincos}{\begin{array}{c} \cos \\ \sin \end{array}\!\!}
\newcommand{\CvBc}[1]{\left\{\:#1\:\right\}}
\newcommand*{\one}{ {{\rm 1\mkern-1.5mu}\!{\rm I} }}
\newcommand{\RR}{\mathbb{R}}
\newcommand{\R}{\mathbb{R}}
\newcommand{\myEnd}{\hfill$\square$}
\newcommand{\ds}{\displaystyle}
\newcommand{\Shi}{\text{Shi}}
\newcommand{\Chi}{\text{Chi}}
\newcommand{\Daw}{\ensuremath{\mathrm{daw} }}
\newcommand{\Erf}{\ensuremath{\mathrm{erf} }}
\newcommand{\Erfi}{\ensuremath{\mathrm{erfi} }}
\newcommand{\Erfc}{\ensuremath{\mathrm{erfc} }}
\newcommand{\He}{\ensuremath{\mathrm{He} }}

\def\crtL{\theta}

\newcommand{\conRef}[1]{(#1)}

\DeclareMathOperator{\Lip}{\mathit{L}}
\DeclareMathOperator{\LIP}{Lip}
\DeclareMathOperator{\lip}{\mathit{l}}

\DeclareMathOperator{\Vip}{\overline{\varsigma}}
\DeclareMathOperator{\vip}{\underline{\varsigma}}
\DeclareMathOperator{\vv}{\varsigma}
% }}}

\begin{document}

\section{Introduction}
\begin{enumerate}
	\item 
In the first paragraph of page 2, replace the sentence "...where the initial condition $u_0$ satisfying..." with the sentence "...where the initial condition $u_0$ \textit{satisfies}..."

	\item 
In the third paragraph of page 2, replace the sentence "...is not included in STF-SHE (1), by using..." with the sentence "...is not included in \textit{the} STF-SHE (1), \textit{then} by using...". 

	\item
In the fourth paragraph of page 2, replace the sentence "...by considering the heat equation in material..." with the sentence "...by considering the heat equation in \textit{materials}..."

	\item
In the first paragraph on page 3, replace the sentence "...of the corresponding STF-SHE..." with the sentence "...of the corresponding \textit{the} STF-SHE...".

	\item 
In the first paragraph on page 3, replace the sentence "While in [19], for STF-SHE..." with the sentence  "While in [19], for \textit{the} STF-SHE...".

	\item
In the first paragraph on page 3, replace the sentence "...of the solutions of this equation grows exponentially..." with the sentence "...of the solutions of this equation \textit{grow} exponentially...".

	\item
In the first paragraph on page 3, replace the sentence "... studied non-linear noise excitation of the STF-SHE (1) in bounded domain..." with the sentence "... studied \textit{the} non-linear noise excitation of the STF-SHE (1) in \textit{a} bounded domain..."

	\item
In the second paragraph of page 3, replace the sentence "...of the parameters $\alpha$, $\beta$ in STF-SHE (1)." with the sentence "...of the parameters $\alpha$, $\beta$ in \textit{the} STF-SHE (1)."

	\item
In the third paragraph on page 3, replace the sentence "...for several classes of SPDE driven..." with the sentence "...for several classes of \textit{SPDEs} driven...".

	\item
In general, "\textit{weakly full intermittent}" should be replaced with "\textit{weakly fully intermittent}". The first occurrence is in paragraph 3 on page 3. 

	\item 
In the third paragraph on page 3, replace the sentence "...based on the estimates of moments..." with the sentence "...based on the estimates of \textit{the} moments..." 

	\item 
In the third paragraph on page 3, replace the sentence "...to STF-SHE (1). The first novelty..." with the sentence  "...to \textit{the} STF-SHE (1). The first novelty..."

	\item
In the third paragraph on page 3, remove the phrase "and so on" in the sentence "...[11], [19] and so on, in which" and replace it with "...[11], [19], in which...".

	\item 
In the third paragraph on page 3, replace the sentence "Secondly, the Green function..." with the sentence "Secondly, the \textit{Green's} function..."

	\item 
In the third paragraph on page 3, replace the sentence "...of the solution to STF-SHE..." with the sentence  "...of the solution to \textit{the} STF-SHE..."

	\item
In the fourth paragraph on page 3, replace the sentence "...and the Green function..." with the sentence  "...and the \textit{Green's} function..."
\end{enumerate}
		
\section{Preliminaries}
\subsection{The spatially inhomogeneous white noise}
\begin{enumerate}
	\item 
How is it true that the following Ito isometry holds:
	\[
		\mathbb{E}\left[ |(u\cdot F_\rho)(t)|^2 \right] = \mathbb{E}\left[ \int_0^t \int_{\R} u^2(s,x) \ud s \rho(\ud x) \right]?
	\]
$(u\cdot F_\rho)(t)$ is a Walsh integral, not an Ito integral. See Proposition 5.18 of the mini course for the case of a step stochastic integral. Also, I can not find a reference anywhere that calculates the second moment of a Walsh integral.
	
	\item 
In (A.2) on page 5, what is the reason for the $u$ subscript in $C_u$? It seems like $C_u$ only depends on $\rho$ and so it would make more sense to denote it as $C_\rho$.

	\item
In Remark 1 on page 5, replace the phrase "the author in [27]" with the phrase "Z\"{a}hle [27]".

	\item 
In (B) on page 6, the constant $K_u$ depends on $\sigma$ and does not depend on $u$, so $K_u$ should be written as $K_\sigma$. 

	\item 
In the second paragraph on page 6 replace the sentence "...and continuous random field $u=$..." with the sentence "...and continuous random \textit{fields} $u=$..."

	\item 
In the third paragraph on page 6 replace the sentence "...for the solution to STF-SHE..." with the sentence "...for the solution to \textit{the} STF-SHE..."

	\item
I just realized that in $C_u$ and $C_l$ and $K_u$ and $K_l$ that the $u$ and $l$ mean upper and lower respectively. However, since the solution is also denoted by $u$, this may cause some confusion because it seems like $C_u$ and $K_u$ depend on the solution $u(t,x)$ but this is not the case. 

\end{enumerate}

\subsection{Some properties of Green function}
\begin{enumerate}
	\item
The title of this subsection should be "Some properties of \textit{the} \textit{Green's} function". 

	\item 
The first sentence of the subsection needs to be reworded and the use of the phrase "i.e." is used incorrectly. I would word the first sentence as follows: 
	
Consider the fundamental solution (or Green's function) $G_t(x)$ of the fractional heat type equation:
	\[
		\partial^{\beta}_t G_t(x) = -\nu(-\Delta)^{\alpha/2}G_t(x).
	\]
Due to our need to apply some properties of $G_t(x)$, we must assume $\nu>0$, $\beta \in (0,1)$ and $\alpha \in (0,2]$. 
	
	\item
Does the Green's function not depend of the fractional integral operator $I_t^{1-\beta}$? Why do we only consider the Greens function for 
	\[
		\partial^{\beta}_t G_t(x) = -\nu(-\Delta)^{\alpha/2}G_t(x)
	\]
and not the equation 
	\[
		\partial^{\beta}_t G_t(x) = -\nu(-\Delta)^{\alpha/2}G_t(x) + I_t^{1-\beta}(\lambda f(t,x))
	\]
for some deterministic $f$?

	\item 
In the first paragraph on page 7, replace the sentence "...the first passage time of an $\beta$-stable..." with the sentence "...the first \textit{passage-time} of \textit{a} $\beta$-stable..."

	\item 
In the first paragraph on page 7, what is meant by the sentence "...or the inverse stable subordinator of index $\beta$."?  
\end{enumerate}


\section{Moments bounds and intermittency}
\begin{enumerate}
	\item
The title should be changed to "\textit{Moment} bounds and intermittency" 
	
	\item
I would remove the word "some" before the phrase "some Banach space $\mathbb{B}_{p,\theta}$. This is because we already defined what the Banach space $\mathbb{B}_{p,\theta}$ is and so it is no longer just "some" Banach space. This error is also made in the first paragraph of the subsection Upper bounds for the $p$-th moment. 
	
	\item
In the first paragraph of this section, change "...the solution $u(t,x)$ to STF-SHE..." to "...the solution $u(t,x)$ to \textit{the} STF-SHE...".

	\item 
In the first paragraph of this section, there is an incomplete thought in the sentence "...we will prove the uniformly lower Lyapunov's exponent of the solution $u(t,x)$ to STF-SHE (1), which implies...". It was never stated what you are going to prove that will imply the solution grows at least exponentially. And in addition, this paper finds the Lyapunov exponent for the moments of $u(t,x)$, not $u(t,x)$ itself so we can only use this to say $\mathbb{E}|u(t,x)|^p$ grows at least exponentially. 

	\item 
In Definition 2 on page 7, why do we use this $G_t(x)$ which is the fundamental solution to equation (3 to solve equation (1)? Equation (1) contains the fractional integral operator while equation (3) does not. It seems that the fundamental solution should also depend on $I^{1-\beta}_t$. 
\end{enumerate}

\subsection{Upper bounds for the $p$-th moment}
\begin{enumerate}
	\item 
In the statement of Theorem 1, change the phrases "... $u_0(x)$ satisfying..." to the phrase "... $u_0(x)$ \textit{satisfies}...".

	\item
I believe there is an error in the proof equation (6). To keep everything in order, I will discuss this later on. 

	\item 
In 1 from Remark 2, replace "...order of fractional operator..." with "...order of \textit{the} fractional operator...".

	\item
In 2 from Remark 2, replace "...$\beta \in (0,1)$, one obtains..." with "...$\beta \in (0,1)$, \textit{then} one obtains...".

	\item 
In the last line of 2 from Remark 2, replace "Especially, the..." with "\textit{However}, the...".

	\item
In 3 from Remark 2, replace "...reduces to the quantities..." with "...\textit{reduce} to the quantities...".
	
	\item
In 3 from Remark 2, replace "...which coincides the results..." with "...which coincides \textit{with} the results...".

	\item 
In 3 from Remark 2, replace "...$u(t,x)$ is less..." with "...$u(t,x)$ \textit{are} less...".

	\item 
In 3 from Remark 2, replace "...the order of Holder regularity turn to be less $\frac{1}{4}$..." with "...the \textit{orders} of Holder regularity \textit{become less than} $\frac{1}{4}$..."

	\item
In 3 from Remark 2, replace "...in studying stochastic heat..." with "...in studying \textit{the} stochastic heat..." 
\end{enumerate}
	
\textbf{Lemma 2.} 
\begin{enumerate}
	\item 
Is $C_3$ independent of $\beta$? In the proof of Lemma 2, we use Lemma 7 which says that when $\alpha > d$ then 
		\[
			G_t(x) < C_2\left(t^{-\beta d/\alpha} \wedge \frac{t^\beta}{|x|^{d+\alpha}} \right).
		\]
	Thus $C_2$ is dependent on $G_t(x)$ and $G_t(x)$ is dependent on $\beta$. The proof of Lemma 2 uses Lemma 9 to come up with $C_3$ and after examining the proof one can see that $C_3$ depends on $C_2$ which depends on $\beta$. In addition it depends on $C_u$ which depends on $\rho$. So it seems that $C_3$ depends on both $\rho$ and $\beta$. 
	
	\item 
I think the last inequality 
	\[
		\le 4p\Norm{u}_{p,\theta}^2 \sup_{(t,x)\in[0,T]\times \R} \int_0^t\int_{\R} G^2_s(x-y)\exp\left\{-\frac{2\theta s}{p} \right\} \rho(\ud y) \ud s.
	\]
should be an equality which can be seen by applying the substitution $u=t-s$ to the previous line. 
	\item
I believe the result is slightly off. I think what we can say is that 
	\begin{equation}\label{L:Lemma2cor}
		\Norm{\mathbb{S}u}_{p,\theta}^2 \le 4pC_3C_{\bar{\delta}}^*(\alpha,\beta,2\theta/p) \Norm{u}_{p,\theta}^2.
	\end{equation}
Notice that the last time of the proof on page 9 gives us that 
	\[
		\Norm{\mathbb{S}u}_{p,\theta}^2 \le 4p\Norm{u}_{p,\theta}^2 \sup_{(t,x)\in[0,T]\times \R} \int_0^t\int_{\R} G^2_s(x-y)\exp\left\{-\frac{2\theta s}{p} \right\} \rho(\ud y) \ud s.
	\]
Therefore to apply Lemma 9 at this step, we need to let $\theta' = 2\theta/p$ where $\theta'$ is the parameter from Lemma 9. Thus we can say that 
	\begin{align*}
		4p\Norm{u}_{p,\theta}^2 \sup_{(t,x)\in[0,T]\times \R} \int_0^t\int_{\R} G^2_s(x-y)\exp\left\{-\frac{2\theta s}{p} \right\} \rho(\ud y) \ud s &\le 4pC_3C_{\bar{\delta}}^*{\alpha,\beta,\theta'}\Norm{u}_{p,\theta}^2
		\\& = 4pC_3C_{\bar{\delta}}^*(\alpha,\beta,2\theta/p) \Norm{u}_{p,\theta}^2.
	\end{align*}
\end{enumerate} 

\textbf{Lemma 3.}
\begin{enumerate}
	\item
Is $C_3'$ not the same time as $C_3$ from Lemma 2? Why can't we apply Lemma 2 directly and say that 
	\begin{align*}
		\Norm{\mathbb{S}u-\mathbb{S}v}_{p,\theta}^2 &= \Norm{\mathbb{S}(u-v)}_{p,\theta}^2
		\\&\le 4pC_3C_{\bar{\delta}}^*(\alpha,\beta,2\theta/p) \Norm{u-v}_{p,\theta}^2
	\end{align*}
where above I added the possible correction of Lemma 2 \eqref{L:Lemma2cor}. 
\end{enumerate}

\textbf{Lemma 11:}

Before I talk about Lemma 4, I will talk about Lemma 11 since the proof of Lemma 4 relies on Lemma 11.

Essentially, what Lemma 11 showed was that
	\[
		\int_0^{t_2}\int_{\R} |G_{t-2-s}(x_2-y) - G_{t_1-s}(x_1-y)|^2 \rho(\ud y) \ud s \le 2(\mathcal{L} + \mathcal{M} + \mathcal{N})
	\]
where 
	\begin{align*}
		\mathcal{L} &=	\int_{t_1}^{t_2}\int_{\R} |G_{t-2-s}(x_2-y)|^2 \rho(\ud y) \ud s \le C|t_2-t_1|^{\frac{\alpha+\beta(\bar{\delta}-2)}{\alpha}}
	\\	\mathcal{M} &= 	\int_0^{t_1}\int_{\R} |G_{t-2-s}(x_1-y) - G_{t_1-s}(x_1-y)|^2 \rho(\ud y) \ud s \le C|t_2-t_1|^{\frac{\alpha+\beta(\bar{\delta}-2)}{\alpha}}
	\\ \mathcal{N} &= 	\int_0^{t_2}\int_{\R} |G_{t-2-s}(x_2-y) - G_{t_2-s}(x_1-y)|^2 \rho(\ud y) \ud s \le C |x_2-x_1|^{2\kappa}, \quad \text{for }0 < \kappa < \frac{\alpha+\beta(\bar{\delta}-2)}{\beta}.
	\end{align*}

I don't see in the proof where $\mu$ comes in to play and the only mention of $\mu$ is in the statement of the Lemma, $\mu$ never appears in the proof. 

I believe what this proof shows was that 
	\[
			\int_0^{t_2}\int_{\R} |G_{t-2-s}(x_2-y) - G_{t_1-s}(x_1-y)|^2 \rho(\ud y) \ud s \le C \left( |t_2-t_1|^{\frac{\alpha+\beta(\bar{\delta}-2)}{\alpha}} + |x_2-x_1|^{2\kappa} \right)
	\]
where 
	\[
		0 < \kappa < \frac{\alpha+\beta(\bar{\delta}-2)}{\beta}.
	\]

\textbf{Lemma 4:}
\begin{enumerate}
\item It seems unclear how Lemma 11 can conclude the proof. To summarize, the last line of the proof of Lemma is says 	
	\begin{align*}
		&\mathbb{E}\left[ |\mathbb{S}u(t_2,x_2) - \mathbb{S}u(t_1,x_1)|^p \right]
		\\& \le 2^{p-1}(4p)^{p/2} \left(\sup_{(s,y)\in[0,T]\times \R} \mathbb{E}\left[|u(s,y)|^p\right]\right) \left( \int_0^{t_1} (G_{t_2-s}(x_2-y)-G_{t_1-s}(x_1-y))^2 \rho(\ud y) \ud s \right)^{p/2}
		\\& \quad \quad+ 2^{p-1}(4p)^{p/2} \left(\sup_{(s,y)\in[0,T]\times \R} \mathbb{E}\left[|u(s,y)|^p\right]\right)\left( \int_{t_1}^{t_2} G^2_{t_2-s}(x_2-y) \rho(\ud y) \ud s \right)^{p/2}
	\end{align*}
and its stated the result follows from Lemma 11. However I do not see how this is the case. What I think we can show is that 
	\begin{align*}
		& \le 2^{p-1}(4p)^{p/2} \left(\sup_{(s,y)\in[0,T]\times \R} \mathbb{E}\left[|u(s,y)|^p\right]\right) \left( \int_0^{t_1} (G_{t_2-s}(x_2-y)-G_{t_1-s}(x_1-y))^2 \rho(\ud y) \ud s \right)^{p/2}
		\\& \quad \quad+ 2^{p-1}(4p)^{p/2} \left(\sup_{(s,y)\in[0,T]\times \R} \mathbb{E}\left[|u(s,y)|^p\right]\right)\left( \int_{t_1}^{t_2} G^2_{t_2-s}(x_2-y) \rho(\ud y) \ud s \right)^{p/2}
		\\& \le 2^{p-1}(4p)^{p/2} \left(\sup_{(s,y)\in[0,T]\times \R} \mathbb{E}\left[|u(s,y)|^p\right]\right) \left( \int_0^{t_2} (G_{t_2-s}(x_2-y)-G_{t_1-s}(x_1-y))^2 \rho(\ud y) \ud s \right)^{p/2}
		\\& \quad \quad+ 2^{p-1}(4p)^{p/2} \left(\sup_{(s,y)\in[0,T]\times \R} \mathbb{E}\left[|u(s,y)|^p\right]\right)\left( \int_{t_1}^{t_2} G^2_{t_2-s}(x_2-y) \rho(\ud y) \ud s \right)^{p/2}
		\\& \le 2^{p-1}(4p)^{p/2} \left(\sup_{(s,y)\in[0,T]\times \R} \mathbb{E}\left[|u(s,y)|^p\right]\right) \left( 2(\mathcal{L}+\mathcal{M} + \mathcal{N}) \right)^{p/2}
		\\& \quad \quad+ 2^{p-1}(4p)^{p/2} \left(\sup_{(s,y)\in[0,T]\times \R} \mathbb{E}\left[|u(s,y)|^p\right]\right)\left( \mathcal{L} \right)^{p/2}
		\\& \le 2^{p-1}(4p)^{p/2} \left(\sup_{(s,y)\in[0,T]\times \R} \mathbb{E}\left[|u(s,y)|^p\right]\right) \left( C\left(|t_2-t_1|^{\frac{\alpha + \beta(\hat{\delta} -2)}{\alpha}} + |x_2-x_1|^{2\kappa} \right)\right)^{p/2}
		\\& \quad \quad+ 2^{p-1}(4p)^{p/2} \left(\sup_{(s,y)\in[0,T]\times \R} \mathbb{E}\left[|u(s,y)|^p\right]\right)\left( C\left(|t_2-t_1|^{\frac{\alpha + \beta(\hat{\delta} -2)}{\alpha}}  \right)^{p/2} \right)
	\end{align*}
where 
	\[
			0 < \kappa < \frac{\alpha+\beta(\bar{\delta}-2)}{\beta}.
	\]
This does not reduce to statement of the Lemma. However, through repeated use of the inequality $(a+b)^p \le 2^{p-1}(a^p+b^p)$ we can say that 
	\begin{align}
		&2^{p-1}(4p)^{p/2} \left(\sup_{(s,y)\in[0,T]\times \R} \mathbb{E}\left[|u(s,y)|^p\right]\right) \left( C\left(|t_2-t_1|^{\frac{\alpha + \beta(\hat{\delta} -2)}{\alpha}} + |x_2-x_1|^{2\kappa} \right)\right)^{p/2} \nonumber
	\\ & \quad \quad+ 2^{p-1}(4p)^{p/2} \left(\sup_{(s,y)\in[0,T]\times \R} \mathbb{E}\left[|u(s,y)|^p\right]\right)\left( C\left(|t_2-t_1|^{\frac{\alpha + \beta(\hat{\delta} -2)}{\alpha}}  \right)^{p/2} \right) \nonumber
	\\ & \le 2^{(3p-4)/2}(4p)^{p/2} C\left(\sup_{(s,y)\in[0,T]\times \R} \mathbb{E}\left[|u(s,y)|^p\right]\right) \left(|t_2-t_1|^{\frac{p\alpha + p\beta(\hat{\delta} -2)}{2\alpha}} + |x_2-x_1|^{p\kappa} \right) \label{E:holdercont}
	\end{align}
and this inequality is essentially just as good as the one stated in the statement of Lemma 4. 

\item I believe there is a problem when citing Kolmogrov's continuity theorem. The statement of the theorem that I am aware of is the following: 

\textit{Kolmogrov Continuity Theorem:} Consider a real valued processes $\{X_t\}_{t \times \R}$. If there exists positive constants $p, b$ and $C$ such that 
		\[
			\mathbb{E}\left[ |X_t-X_s|^p \right] \le C |t-s|^{1+b}
		\]
then there is a modification of $X$ that is locally $\gamma$-Holder continuous for $0<\gamma < b/p$. 

Thus if we apply this theorem with \eqref{E:holdercont} then what we see is that there is a modification of the solution that is $\mu_1$-Holder continuous in time and $\mu_2$-Holder continuous in space where 
	\[
		0 < \mu_1 <\left( \frac{p\alpha + p\beta(\hat{\delta} -2)}{2\alpha} - 1 \right)/p
	\]
and 
	\[
		0 < \mu_2 < (p\kappa -1)/p.
	\]
\end{enumerate}

\textbf{Proof of Theorem 1 (Part 1):} 
\begin{enumerate}
	\item 
Replace the first sentence, "Let us first..." with the following sentence: "We start by proving part (I).". 
	
	\item 
In the second sentence of the proof, replace the phase "...we can easily get the first term in the right hand of [5]..." with the phrase "...we can easily \textit{see} that the first term \textit{on} the right hand \textit{side} of [5],...".
	
	\item
In the fourth line of the proof, what "contractivity" of $\mathbb{T}_{0,t}u_0(x)$ are you referring too? 

	\item 
How is the first inequality in the set of centered equations on page 11 shown? It particular the inequality 
	\[
		\sup_{(t,x)\in[0,T]\times \R} \mathbb{E}\left[ e^{-\theta t} \left|\int_{\R} G_t(x-y)u_0(y) \ud y \right|^p \right] \le \sup_{(t,x)\in[0,T]\times \R} \mathbb{E}\left[ e^{-\theta t} \int_{\R} G_t(x-y)|u_0(y)|^p \ud y  \right].
	\]
	 
	\item
How is the second inequality in the set of centered equations on page 11 shown? It particular the inequality
	\[
		\sup_{(t,x)\in[0,T]\times \R} \mathbb{E}\left[ e^{\theta t} \int_{\R} G_t(x-y)|u_0(y)|^p \ud y \right]  \le \sup_{(t,x)\in[0,T]\times \R} \mathbb{E}\left[|u_0(y)|^p\right]e^{-\theta t}.
	\]
What happened to 
\[
\int_{\R}G_t(x-y) ]\ud y ?
\]
Does this integral equal $1$ for all $x \in \R$ and $t>0$? I not, then is this integral uniformly bounded for all $x \in \R$ and $t>0$? 

	\item 
In the second paragraph on page 12, why are we now assuming that $\hat{\delta} \in (0,1)$? Assumption (A.2) says that $\hat{\delta} \in (0,2]$. 

	\item 
In the second paragraph on page 12, replace the phrase "...it is known that the mapping..." with the phrase "...\textit{we see} that the mapping..."
\end{enumerate}

\textbf{Proof of Theorem 1 (Part 2):}
\begin{enumerate}
	\item 
In the third paragraph on page 12 replace the sentence "From now on, we will formulate the proof of the second point." with the sentence "We now formulate the proof of part 2." 

	\item 
In the third paragraph on page 12, why is $\hat{\delta} \in (0,1]$? This is different from Assumption A.2.

	\item 
In the third paragraph on page 12, how does the constant $C$ only depend on $C_3$ and $C_3'$? The larger $K_u$ and $p$ become then the larger we need to choose $C$ so it seems that $C$ depends on $K_u$ and $p$. 

	\item 
In the third paragraph on page 12, what is the motivation for choosing that specific $\theta$, namely, 
	\[
		\theta = C K_u^{\frac{2\alpha}{\alpha + \beta(\hat{\delta}-2)}} p^{1+\frac{\alpha}{\alpha + \beta(\hat{\delta}-2)}}.
	\]
Why not just pick an arbitrary $\theta > \theta_0$ and continue the proof this way which would lead to the expression that for any fixed $\theta>\theta_0$, there exists a $C_1>0$ such that
	\[
		C_1^{p/2}\left( 1 + \sup_{x\in \R} \mathbb{E}\left[|u_0(x)|^p\right] \right) \exp(\theta t).
	\]
	
	\item
In the third paragraph on page 12, what does it mean that for the particular $\theta$ that $pC_{\hat{\delta}}^*(\alpha, \beta, \theta)$ is bounded? Does this mean that it is finite? 

	\item
In the third paragraph on page 12, replace "...to STF-SHE.." with "...to \textit{the} STF-SHE..."
	
	\item 
In the third paragraph on page 12, replace the phrase "...into $\mathbb{B}_{p,\theta}$, we know that..." with the phrase "...into $\mathbb{B}_{p,\theta}$, \textit{so} we know that..."

	\item 
There is a mistake on the first centered equation on the top of page 13. We have that 	
	\[
		\left( \mathbb{E}\left[ e^{-\theta t} |u(t,x)|^p \right] \right)^{2/p} \le \Norm{u}_{\theta,p}^2 \le 2\sup_{x \in \R} \mathbb{E}\left[ |u_0(x)|^p \right]^{p/2} + 2C'.
	\]
and this implies that
	\[
			\mathbb{E}\left[ |u(t,x)|^p \right]  \le \left(2\sup_{x \in \R} \mathbb{E}\left[ |u_0(x)|^p \right]^{p/2} + 2C'\right)^{p/2}e^{\theta t }.
	\]
From this we can say that 
	\[
			\mathbb{E}\left[ |u(t,x)|^p \right]  \le 2^{(p-2)/2}\tilde{C}^{p/2} \left(1 + \mathbb{E}\left[ |u_0(x)|^p \right]\right) e^{\theta t}
	\]
where 
	\[
		\tilde{C} = \max\{ 2, 2C' \}.
	\]
\end{enumerate}

\textbf{Proof of Theorem 1 (Part 2):}
\begin{enumerate}
	\item
Does the $C$ appearing in the third set of centered equations on page 13 not depend on $t_1$ and $t_2$? In particular the $C$ from 
	\begin{equation}\label{E:timeinc}
		\sup_{y\in\R} \mathbb{E}\left[|u_0(y)|^p\right]|t_2-t_1|^p \left| \int_{\R} \frac{\partial}{\partial_t} G_t(x-y) \vert_{t=\xi} \ud y \right|^p \le C|t_2-t_1|^p.
	\end{equation} 
Because if we use the bound 
	\[
		\frac{\partial}{\partial t}G_t(x-y) \le \frac{C}{t}G_t(x-y)
	\]
then we get
	\begin{align*}
		\int_{\R} \frac{\partial}{\partial_t} G_t(x-y) \vert_{t=\xi} \ud y  &\le \int_{\R}  G_\xi(x-y) \ud y.
	\end{align*}
It is claimed that 
	\[
	\left|\int_{\R}  G_t(x-y) \ud y \right|
	\]
is bounded but what does that mean? Is the bound uniform in $x$ and $t$? If the bound is not uniform in $t$ then the $C$ from \eqref{E:timeinc} depends on $t_1$ and $t_2$ which cant happen. 

	\item
In the last paragraph on page 13, replace the phrase "We will therefore..." with "We will \textit{now}...".

	\item
I believe there is something missing from the first inequality on the top of page 14. What is the $\mathbb{E}$ acting on? Also, shouldn't is say $\rho(\ud y) \ud s$? Here is the inequality in question:
	\[
		\mathbb{E}\left[ \left| \int_{t_1}^{t_2} \int_{\R} G_{t_2-s}(x-y)\sigma(u(s,y))F_{\rho}(\ud s, \ud y) \right|^p \right] \le C \mathbb{E} \left[ \left| \int_{t_1}^{t_2} \int_{\R} G^2_{t_2-s}(x-y)\rho(\ud s) \ud y \right|^{p/2} \right]
	\]
	
	\item 
Let's consider the very next inequality,
	\begin{align*}
		&C \mathbb{E} \left[ \left| \int_{t_1}^{t_2} \int_{\R} G^2_{t_2-s}(x-y)\rho(\ud s) \ud y \right|^{p/2} \right] 
	\\ & \le C \int_{t_1}{t_2} (t_2 - s)^{-2\beta/\alpha} \int_{\R}1_{\left\{ |x-y| \le (t_2-s)^{\beta / \alpha} \right\}}\rho(\ud y) \ud s 
	\\ & \quad \quad + C \int_{t_1}{t_2} (t_2-s)^{2\beta} \int_{\R} \frac{1}{|x-y|^{2(1+\alpha)}}1_{\left\{ |x-y| > (t_2-s)^{\beta / \alpha} \right\}} \rho(\ud y) \ud s.
	\end{align*}
What happened to the power $p/2$ in the first integral? The entire integral on the right hand side should be raised to the power $p/2$.

	\item
Now lets consider the next inequality from the top of page 14 with the correction mentioned above,
	\begin{align*}
		& C \bigg(\int_{t_1}{t_2} (t_2 - s)^{-2\beta/\alpha} \int_{\R}1_{\left\{ |x-y| \le (t_2-s)^{\beta / \alpha} \right\}}\rho(\ud y) \ud s 
	\\ & \quad \quad + C \int_{t_1}{t_2} (t_2-s)^{2\beta} \int_{\R} \frac{1}{|x-y|^{2(1+\alpha)}}1_{\left\{ |x-y| > (t_2-s)^{\beta / \alpha} \right\}} \rho(\ud y) \ud s\bigg)^{p/2}
	\\ & \le C(t_2-t_1)^{p\mu}. 		
	\end{align*}
First, where is this $\mu$ coming from? It seems like that all we can say is that 
	\begin{align*}
		& C \bigg(\int_{t_1}{t_2} (t_2 - s)^{-2\beta/\alpha} \int_{\R}1_{\left\{ |x-y| \le (t_2-s)^{\beta / \alpha} \right\}}\rho(\ud y) \ud s 
		\\ & \quad \quad + C \int_{t_1}{t_2} (t_2-s)^{2\beta} \int_{\R} \frac{1}{|x-y|^{2(1+\alpha)}}1_{\left\{ |x-y| > (t_2-s)^{\beta / \alpha} \right\}} \rho(\ud y) \ud s\bigg)^{p/2}
		\\ & C\left[ \left( t_2-t_1\right)^{\frac{\alpha+\beta(\bar{\delta}-2)}{\alpha}} \right]^{p/2}
		\\& = C\left( t_2-t_1\right)^{\frac{p\alpha+p\beta(\bar{\delta}-2)}{2\alpha}} 
	\end{align*}	
where above we used the bound from $\mathcal{L}$ from Lemma 11. 

	\item
On page 13, replace the phrase "Secondly, let us..." with "\textit{We now} consider...".

	\item
What happened to the $C|x_2-x_1|^{p\varrho}$ term in the upper bound given in the first centered equation on page 15. Following what is written in this proof we have that 
	\begin{align*}
		&\mathbb{E}\left[ |u(t,x_2)-u(t,x_1)|^p \right] 
	\\ &\le 2^{p-1} \Bigg( \mathbb{E}\left[ \left| \int_{\R} [G_t(x_2-y)-G_t(x_1-y)]u_0(y) \ud y \right|^p \right]
	\\ & +  \mathbb{E}\left[ \left| \int_0^t\int_{\R} [ G_{t-s}(x_2-y)-G_{t-s}(x_2-y) ] \sigma(u(s,y)) F_{\rho}(\ud s, \ud y)  \right|^p \right] \Bigg)
	\end{align*}
and in the proof it is claimed that 
	\[
		 \mathbb{E}\left[ \left| \int_{\R} [G_t(x_2-y)-G_t(x_1-y)]u_0(y) \ud y \right|^p \right] \le C|x_2-x_1|^{p\varrho}
	\]
and that 
	\[
		\mathbb{E}\left[ \left| \int_0^t\int_{\R} [ G_{t-s}(x_2-y)-G_{t-s}(x_2-y) ] \sigma(u(s,y)) F_{\rho}(\ud s, \ud y)  \right|^p \right]  \le C|x_2-x_1|^{p\kappa}
	\]
and thus we must have 
	\[
		\mathbb{E}\left[ |u(t,x_2)-u(t,x_1)|^p \right]  \le 2^{p-1}C\left(|x_2-x_1|^{p\varrho}+ |x_2-x_1|^{p\kappa} \right).
	\]
\end{enumerate}

\subsection{Lower bound for the second moment}
\begin{enumerate}
	\item 
The first sentence is unclear. There needs to be a mention of "what" is implying the exponential growth of the solution. All that is said that "...we are going to consider the uniformly lower Lyapunov exponent $\underline{\mathcal{L}}(p,x), \; p \ge 2$ of the solution $u(t,x)$ to the STF-SHE, which implies...". However we need to mention what about the lower Lyapunov exponent that implies the exponential growth. Otherwise the sentence is vague and unclear.
	
	\item 
The second sentence needs to be reworded. Here is a possible way to do so: "In fact, under assumption C, if we can show that $\int_{x \in \R}\underline{\mathcal{L}}(2,x) >0$ then the solution $u(t,x)$ will exhibit the weakly full intermittency property. 

	\item
Remove the phrase "Then we can state from the last sentence of the second paragraph. Just say "The main theorem in this subsection is as follows.".
\end{enumerate}

\textbf{Theorem 2}
\begin{enumerate}
	\item 
In the statement of part 2 on the top of page 16, what is the point of $T$. All it says to fix $T$ large and $T>1$ and then there is never a mention of what this $T$ is for since equation 17 never mentions anything about $T$.
\end{enumerate}

\textbf{Proof of Theorem 2 (Part 1):}
\begin{enumerate}
	\item 
The word "case" is used wrong. Instead of saying case 1 or case 2, you should say part 1 and part 2. 

	\item 
In the very first line of the proof, it says "...by the property of the Green function $G_t(x)$...". What property are you referring too? In addition, $G_t(x)$ should be referred to as the "Green's function", not the "Green function". 

	\item 
There is a claim made in this paper as well as the paper \cite{FooNane2021} that by the Walsh isometry then we have that 
	\[
		\mathbb{E}\left[ \left| \int_0^t\int_{\R} G_{t-s}(x-y)\sigma(u(s,y)) F_{\rho}(\ud s \ud y) \right|^2\right] = \int_0^t\int_{\R} G^2_{t-s}(x-y)\mathbb{E}\left[|\sigma(u(s,y))|^2\right] \ud s \ud y 
	\]
but this is not a result stated in either Walsh \cite{Walsh} or \cite{Minicourse09}. As far as I know, all that can be said is that 
	\begin{equation}\label{MiniTh5.26}
		\mathbb{E}\left[ \left| \int_0^t\int_{\R} G_{t-s}(x-y)\sigma(u(s,y)) F_{\rho}(\ud s \ud y) \right|^2\right] \le  \int_0^t\int_{\R} \prod_{i=1}^2G^2_{t-s}(x-x_i)\mathbb{E}\left[|\sigma(u(s,x_i))|^2\right] K(\ud s \; \ud x_1 \; \ud x_2 )
	\end{equation}
where $K$ is the dominating measure for $F_{\rho}$. See for example, \cite{Minicourse09} Theorem 5.26. If I am correct then the rest of the proof in invalid from this step. In order for the proof to work, equality is necessary in \eqref{MiniTh5.26}. I will continue reviewing the proof ignoring this concern. 

	\item 
Change the phrase "Now let us assume that $T_0=$..." to the phrase "Now define $T_0:=$...". 
	
	\item 
There is something wrong with the definition of $T_0$ as it is stated. For any fixed $t\in (0,T]$, what does the following mean:
	\[
		\min\left\{ (t-s)^{-\beta/\alpha}, (t-s)^{\frac{\beta}{\alpha}(\underline{\delta}-2)} \right\} \ge C_5K_l \Gamma\left(1 + \frac{\beta(\underline{\delta}-2)}{\alpha}\right)?
	\]
The left hand side is a function of $s$ and the right hand side is a constant. Do you mean that the left hand side is bounded below uniformly in $s$? Did you mean to write something like 
	\[
		\min_{s \in (0,t)}\left\{ (t-s)^{-\beta/\alpha}, (t-s)^{\frac{\beta}{\alpha}(\underline{\delta}-2)} \right\} \ge C_5K_l \Gamma\left(1 + \frac{\beta(\underline{\delta}-2)}{\alpha}\right)?
	\]
In addition, if we examine Lemma 6, it seems like we would want $T_0$ to be something like the following:
	\[
		\inf_{t\in (0,T]} \left\{\min_{s \in (0,t)}\left\{ (t-s)^{-\beta/\alpha}, (t-s)^{\frac{\beta}{\alpha}(\underline{\delta}-2)} \right\} \ge C_5K_l \left(\Gamma\left(1 + \frac{\beta(\underline{\delta}-2)}{\alpha}\right) \right)^{-1} \right\}
	\]
Also, in either case, why does $T_0$ even exit? What if $T_0 = \infty$?  

	\item
If $T_0 = \infty$, ie when 
	\[
		\left\{t \in (0,T] \; : \; \min_{s \in (0,t)}\left\{ (t-s)^{-\beta/\alpha}, (t-s)^{\frac{\beta}{\alpha}(\underline{\delta}-2)} \right\} \ge C_5K_l \left(\Gamma\left(1 + \frac{\beta(\underline{\delta}-2)}{\alpha}\right) \right)^{-1} \right\} = \emptyset
	\]
then the second to last centered equation on the bottom of page 16 is not true.

	\item
I believe the last centered equation on the bottom of page 16 needs some more justification. Its unclear how this follows from Lemma 6. 
\end{enumerate}

\textbf{Theorem 2 (Part 2):}
\begin{enumerate}
	\item 
Label this as "Part 2" not "Case II". 

	\item
 Reword the first sentence of the proof as follows: "We will not prove the second part of Theorem 2.".
 
 	\item 
 Remove the phrase "and etc" from the second sentence of this proof.
 
 	\item 
 Replace the phrase "...decays no faster than polynomial." with the phrase "...has at most polynomial decay."
 
 	\item 
 On page 18, replace the phrase "..."which appearing in the above..." with the phrase "...which appears above in...".
 
 	\item 
 On page 18, remove the phrase "it is led to" and just write "For each $n\in\mathbb{N}$,...".
 
 	\item 
 On page 18, it is said that "We will further decrease the domain on the function $\prod_{i=1}^n G_{s_{i-1}-s_i}^2(y_{i-1}-y_i)$ which has the required lower bound.". What is the "required lower bound" mentioned here? As is, this sentence vague and needs some further explanation. 
 
 	\item
Why is the following true: 
	\[
 	\int_{\mathcal{A}_i} G_{s_i}^2(y_{i-1}-y_i)\rho(\ud y_i) \ge C \min\left\{ s_i^{-\beta/\alpha}, s_i^{\beta(\underline{\delta}-2)/\alpha} \right\}?
	\]

	\item 
How is equation (22) correct? All of the space variable integrals are $\int_{\R}$ while the bound being used is only for $\int_{B(0,(t+t_0)^{\beta/\alpha})}$. 
\end{enumerate}

\addcontentsline{toc}{section}{Bibliography}
\begin{thebibliography}{999}
	\bibitem{Nueman2014}.
	Neuman, E. (2014).
	\newblock Pathwise uniqueness of the Stochastic Heat Equation with Spatially Inhomogeneous White Noise. 
	\newblock {https://arxiv.org/pdf/1403.4491.pdf}

	\bibitem{FooNane2021}.
	Foondun, M., Erkan, N.  (2021).
	\newblock Asymptotic properties of some space-time
	fractional stochastic equations. 
	\newblock {https://arxiv.org/pdf/1505.04615.pdf}
	
	\bibitem{Minicourse09}
	Dalang, R., Khoshnevisan, D., Mueller, C., Nualart, D., and Xiao, Y. (2009).
	\newblock {\it A minicourse on stochastic partial differential equations}.
	\newblock Springer-Verlag, Berlin, 2009. xii+216 pp.
	
	\bibitem[Wal86]{Walsh}
	Walsh, John B.
	\newblock {\it An Introduction to Stochastic Partial Differential Equations}.
	\newblock In: \`Ecole d'\`et\'e de probabilit\'es de
	Saint-Flour, XIV---1984, 265--439.
	Lecture Notes in Math.\ 1180, Springer, Berlin, 1986.

\end{thebibliography}

\end{document}